\documentclass[11pt]{article}

\usepackage[margin=1in]{geometry}
\usepackage{amsmath,amssymb,amsthm}

\newtheorem*{theorem}{Result}

\newcommand{\R}{\mathbb R}
\newcommand{\Z}{\mathbb Z}
\newcommand{\ch}{\operatorname{Ch}}

\title{A Dihedrally Equivariant Ordering of Type \(A\) Chambers}
\author{}
\date{}

\begin{document}
\maketitle

Let \(W=S_n\) be the Weyl group of type \(A_{n-1}\), acting on
\[
  V=\R^n/\langle(1,\ldots,1)\rangle
\]
by permuting coordinates. Let
\[
  \mathcal A_{n-1}
  =
  \{\,x_i=x_j : 1\leq i<j\leq n\,\}
\]
be the Coxeter arrangement of type \(A_{n-1}\). Its chambers are naturally
indexed by the elements of \(W\), equivalently by the \(n!\) linear orders of
the coordinates \(x_1,\ldots,x_n\).

Place the indices \(1,\ldots,n\) on the vertices of a regular \(n\)-gon. This
chooses a dihedral subgroup
\[
  D_n=\langle r,s\rangle \subset S_n,
\]
where
\[
  r:1\mapsto 2\mapsto\cdots\mapsto n\mapsto 1
\]
is a Coxeter element of type \(A_{n-1}\), and \(s\) is a reflection of the
\(n\)-gon. The group \(D_n\) acts on the arrangement, hence on its chambers.

\begin{theorem}
For every \(n\geq 5\), there is a recursively constructed cyclic ordering of all
chambers of the type \(A_{n-1}\) Coxeter arrangement,
\[
  \Phi_n:\ch(\mathcal A_{n-1})\longrightarrow \Z/n!\Z,
\]
with the following dihedral equivariance property.

For every chamber \(C\),
\[
  \Phi_n(r\cdot C)
  =
  \Phi_n(C)+(n-1)!
  \pmod{n!}.
\]
Thus the Coxeter element \(r\) acts on the ordering as a rotation of the cycle
of \(n!\) chambers.

Moreover, the reflection \(s\) acts as a reflection of the same cycle: there is
a constant \(a\in\Z/n!\Z\), depending on the choice of origin of the ordering,
such that
\[
  \Phi_n(s\cdot C)
  =
  a-\Phi_n(C)
  \pmod{n!}.
\]
\end{theorem}

Equivalently, the recursive construction realizes the natural dihedral action
\[
  D_n\curvearrowright \ch(\mathcal A_{n-1})
\]
as the ordinary dihedral symmetry of the cyclically ordered set
\(\Z/n!\Z\). The ordering is produced recursively from rank \(n-2\) to rank
\(n-1\), rather than by choosing representatives of \(D_n\)-orbits.

\end{document}
